\documentclass[sigconf,review,anonymous]{acmart}

\usepackage[ruled,vlined]{algorithm2e}
\usepackage{booktabs}
\usepackage{pgfplots}
\usepackage{float}
\usepackage{tikz}
\usepackage{enumitem}
\usepackage{placeins}
\pgfplotsset{compat=1.18}
\setlength{\textfloatsep}{6pt}
\setlength{\floatsep}{6pt}

\settopmatter{printfolios=true}

\title{PlanGraph: Planning-Feedback Multi-Agent Reasoning for Graph Problems}

\begin{abstract}

Large language models (LLMs) have recently shown strong capabilities in natural language understanding, complex reasoning, and program generation, creating new opportunities for automated graph reasoning. However, graph reasoning involves strong topological dependencies and explicit task constraints. Existing methods often serialize graphs into text and rely on end-to-end reasoning, which makes it difficult to preserve graph topology and satisfy task constraints under long inputs. Another line of work improves code generation with curated task-code repositories or specialized fine-tuning, but such experience-driven solutions are costly to maintain and often generalize poorly to unseen graph tasks.

To address these limitations, we propose \textbf{PlanGraph}, a planning-feedback multi-agent framework for graph reasoning. PlanGraph organizes problem solving into three stages: \textit{planning generation}, \textit{plan-guided retrieval}, and \textit{execution-verified feedback reasoning}. First, question and planning agents transform the natural language task into structured pseudocode that explicitly represents task goals, algorithmic steps, and constraints. Next, a retrieval agent uses the generated plan to retrieve relevant APIs and implementation knowledge from external graph algorithm documentation, and performs snippet refinement or supplementary retrieval when needed. Finally, a coding agent generates executable programs conditioned on both the plan and retrieved knowledge, while an execution-based feedback loop performs testing, error diagnosis, bounded repair, and regeneration. In this way, planning constrains downstream knowledge use and program generation, while feedback corrects implementation deviations and omitted constraints.

We evaluate PlanGraph on five graph reasoning benchmarks: GraphArena, NLGraph, GraphWiz, LLM4DyG, and GraphInstruct. PlanGraph achieves the best performance on four of the five benchmarks, reaches an average exact match score of 97.66\%, and outperforms existing methods overall. The results show that integrating explicit planning, external knowledge retrieval, and program verification into a feedback loop substantially improves the reliability, robustness, and generalization of LLM-based graph reasoning.

\end{abstract}

\keywords{graph reasoning, large language models, multi-agent systems, planning feedback, program verification}

\begin{document}

\maketitle

\section{Introduction}

Graphs are a fundamental data structure for representing complex relational systems in a wide range of real-world applications, including transportation networks, social interactions, biological systems, and software dependency graphs. Many important computational problems can be naturally formulated as graph reasoning tasks, such as shortest path computation, connectivity analysis, graph matching, and combinatorial optimization. Classical solutions to these problems rely on well-established graph algorithms and specialized computation libraries \cite{cormen2009clrs,diestel2017graph,korte2018combinatorial,hagberg2008networkx}, which provide precise and efficient procedures for performing graph analysis. As illustrated in Figure~\ref{fig:task_overview}, graph reasoning tasks typically start from a natural language query describing a graph problem, which must be interpreted over an underlying graph structure before performing algorithmic reasoning.

Recent advances in large language models (LLMs) have significantly expanded the ability of AI systems to perform complex reasoning and program generation tasks. Techniques such as Chain-of-Thought prompting \cite{wei2022chain}, self-consistency decoding \cite{wang2023selfconsistency}, and structured reasoning paradigms including Tree-of-Thoughts \cite{yao2023tree} and Graph-of-Thoughts \cite{besta2024graph} have demonstrated remarkable success in solving mathematical reasoning, logical inference, and algorithmic problem-solving tasks. Furthermore, recent work has shown that LLMs can interact with external tools and generate executable programs to solve complex computational problems \cite{gao2023pal,schick2023toolformer,critic2023}. These developments suggest that LLMs have the potential to serve as general reasoning engines capable of solving structured computational tasks.

However, applying LLMs to graph reasoning problems remains challenging. Unlike natural language or sequential data, graphs contain complex topological structures and relational dependencies that are difficult to represent in textual form. In practice, most existing approaches serialize graph structures into sequences of textual descriptions containing nodes, edges, and attributes. As the size and complexity of the graph increase, such serialized representations become extremely long and difficult for LLMs to interpret. As a result, models often struggle to accurately capture graph topology and relational constraints, leading to unstable reasoning behavior and hallucinated results.

To better understand these challenges, several recent benchmarks have been proposed to evaluate LLM capabilities on graph reasoning tasks. Datasets such as GraphArena \cite{tang2025grapharena}, NLGraph \cite{chen2024nlgraph}, GraphWiz \cite{gong2024graphwiz}, and GraphInstruct \cite{graphinstruct2024} provide systematic evaluations of LLM performance on classical graph computation tasks, including connectivity queries, shortest path search, graph traversal, and combinatorial optimization problems. These studies reveal that although LLMs demonstrate promising reasoning abilities on small or structured examples, their performance often degrades significantly as graph complexity increases.

Recent research has explored several directions for improving LLM-based graph reasoning. One line of work attempts to enhance graph understanding through specialized prompting strategies or instruction tuning \cite{tang2024graphgpt,pan2024graphllm}. Another direction focuses on leveraging code generation capabilities of LLMs to solve graph problems by invoking external graph computation libraries \cite{zhang2025gcoder,pseudocode2025}. In addition, multi-agent frameworks have been proposed to improve reasoning reliability by decomposing tasks into collaborative submodules \cite{li2024graphteam}. These approaches have demonstrated encouraging improvements over direct prompting methods.

\begin{figure}[t]
\centering
\includegraphics[width=\columnwidth]{intro.pdf}
\caption{Illustration of graph reasoning tasks studied in this work. 
A natural language query describing a graph problem is first mapped to a graph structure, 
followed by graph reasoning tasks such as shortest path, maximum clique, and TSP, 
which ultimately produce the final answer.}
\label{fig:task_overview}
\end{figure}

Despite these advances, existing approaches still face two persistent limitations. First, end-to-end reasoning over serialized graph inputs forces the model to jointly handle task interpretation, structural recovery, algorithm selection, and answer generation, which often leads to topological omissions, constraint violations, and unstable outputs as graph complexity grows. Second, many code-generation-based solutions depend on curated task-code repositories, historical exemplars, or specialized fine-tuning. Although such experience-driven strategies can improve performance on seen tasks, they introduce high maintenance costs and do not transfer reliably to new graph reasoning problems.

Inspired by how human experts solve constrained problems through explicit planning and subsequent correction, we propose a planning-feedback reasoning paradigm for graph problem solving. Instead of treating graph reasoning as a single generation step, our approach first constructs a structured plan, then uses that plan to organize knowledge retrieval and program synthesis, and finally uses execution feedback to repair implementation errors or missing constraints.

In this paper, we introduce \textbf{PlanGraph}, a planning-feedback multi-agent framework designed to improve the reliability and generalization of LLM-based graph reasoning systems. PlanGraph decomposes graph problem solving into three collaborative stages: planning generation, plan-guided retrieval, and execution-verified feedback reasoning. First, a planning module converts the natural language task description into structured pseudocode representing task goals, algorithmic steps, and constraints. Second, a retrieval module uses these planning results to retrieve relevant algorithmic knowledge from external sources such as graph library documentation, and refines the retrieved knowledge when coverage is insufficient. Finally, a reasoning module generates executable programs, tests them through execution, and uses bounded repair or regeneration when failures occur.

By explicitly separating planning from execution and coupling generation with feedback, PlanGraph reduces the reasoning burden placed on LLMs and provides a structured pathway for solving complex graph problems. The plan-guided retrieval mechanism enables the system to dynamically access external algorithmic knowledge, while the execution-feedback loop prevents many implementation errors from propagating directly to the final answer.

We evaluate PlanGraph on multiple graph reasoning benchmarks, including GraphArena, NLGraph, GraphWiz, LLM4DyG, and GraphInstruct. Experimental results show that PlanGraph achieves the best performance on four of the five benchmarks and improves overall performance over existing methods.

The main contributions of this work are summarized as follows:

\begin{itemize}[topsep=2pt,itemsep=2pt]

\item We propose \textbf{PlanGraph}, a planning-feedback multi-agent framework for solving graph reasoning tasks with large language models.

\item We introduce a \textbf{plan-guided retrieval mechanism} that leverages structured pseudocode and constraint-aware refinement to improve the precision of external knowledge retrieval.

\item We design an \textbf{execution-verified feedback reasoning pipeline} that integrates program generation, execution testing, bounded repair, and regeneration to improve solution reliability.

\item Extensive experiments on multiple graph reasoning benchmarks demonstrate that PlanGraph achieves the best performance on four of the five benchmarks and improves the robustness of LLM-based graph reasoning systems.

\end{itemize}

\section{Related Work}

\subsection{LLM Reasoning and Tool-Augmented Problem Solving}

Recent advances in large language models (LLMs) have significantly improved their ability to perform complex reasoning tasks. 
Chain-of-Thought (CoT) prompting \cite{wei2022chain} demonstrated that generating intermediate reasoning steps can substantially improve model performance on arithmetic and logical problems. 
Subsequent work proposed decoding strategies such as self-consistency \cite{wang2023selfconsistency}, which aggregates multiple reasoning paths to improve answer reliability. 
More structured reasoning paradigms, including Tree-of-Thoughts \cite{yao2023tree} and Graph-of-Thoughts \cite{besta2024graph}, further extend this idea by enabling language models to explore multiple reasoning branches and structured reasoning paths during problem solving. 
Other prompting strategies such as least-to-most prompting \cite{zhou2023least} decompose complex tasks into a sequence of simpler subproblems, allowing models to solve challenging tasks incrementally.

Beyond prompting-based reasoning, recent research has explored integrating LLMs with external tools and program execution environments. 
Program-Aided Language Models (PAL) \cite{gao2023pal} enable language models to generate executable programs for solving mathematical and symbolic tasks. 
Toolformer \cite{schick2023toolformer} and related systems \cite{critic2023,qin2023toolllm} further extend this paradigm by allowing LLMs to dynamically invoke external APIs during reasoning. 
These tool-augmented frameworks improve reliability by delegating precise computations to external programs or tools. 
However, most existing approaches primarily focus on numerical, mathematical, or symbolic reasoning tasks, and their effectiveness on graph computation problems remains limited, especially when complex graph structures and topological dependencies must be interpreted from textual representations.

\subsection{LLM-based Graph Reasoning}

Recently, several studies have explored applying LLMs to graph analysis and graph reasoning tasks. 
Benchmarks such as GraphArena \cite{tang2025grapharena}, NLGraph \cite{chen2024nlgraph}, GraphWiz \cite{gong2024graphwiz}, GraphInstruct \cite{graphinstruct2024}, and LLM4DyG \cite{llm4dyg2024} evaluate the capability of LLMs to solve classical graph computation problems, including connectivity queries, shortest path search, graph traversal, and combinatorial optimization tasks. 
Experimental results on these benchmarks reveal that although LLMs demonstrate promising reasoning abilities on small or structured graphs, their performance often degrades significantly as graph size and structural complexity increase.

To improve LLM performance on graph reasoning tasks, several research directions have been explored. 
Some studies enhance structural understanding through graph-aware prompting strategies or instruction tuning techniques \cite{tang2024graphgpt,pan2024graphllm}. 
Other approaches leverage the code generation capabilities of LLMs to solve graph problems by invoking external graph computation libraries such as NetworkX \cite{zhang2025gcoder,pseudocode2025}. 
In addition, multi-agent systems such as GraphTeam \cite{li2024graphteam} attempt to improve reasoning reliability by decomposing graph reasoning tasks into collaborative subtasks handled by multiple agents.

Despite these advances, existing approaches still face two major limitations. 
First, many methods rely on end-to-end reasoning over serialized graph representations, which makes it difficult for language models to accurately capture graph topology and relational constraints. 
Second, code-generation-based approaches often rely heavily on curated experience repositories such as task-code pairs or predefined algorithm templates. 
Such experience-driven solutions incur high construction costs and may generalize poorly to unseen graph reasoning tasks. 
Compared with multi-agent collaboration methods such as GraphTeam and program-generation methods such as GCoder, our proposed PlanGraph framework places an explicit planning layer between task understanding and implementation, and further uses execution feedback to revise programs, retrieval targets, or missing constraints when failures occur. This planning-feedback coupling is the main difference between PlanGraph and prior graph reasoning systems.


\section{Methodology}

\begin{figure*}[t]
\centering
\includegraphics[width=0.78\textwidth]{framework.pdf}
\caption{Overview of the PlanGraph framework.}
\label{fig:framework}
\end{figure*}

In this section, we present \textbf{PlanGraph}, a planning-feedback
multi-agent framework designed for solving graph reasoning
problems. Figure~\ref{fig:framework} illustrates the overall
architecture of the proposed system. The key idea of PlanGraph
is to decompose the reasoning process into three stages:
planning generation, plan-guided retrieval, and
execution-verified feedback reasoning.

Instead of relying on direct end-to-end reasoning, PlanGraph
introduces an explicit intermediate planning stage that
structures the reasoning process before code generation,
and an execution-feedback loop that revises incorrect
implementations before they are accepted as final answers.
This design allows the system to better integrate algorithmic
knowledge and external graph computation tools, thereby
improving the reliability and stability of graph reasoning.

\subsection{Framework Overview}

\paragraph{Problem Definition}

Given a graph reasoning task described in natural language
$Q$ and an associated graph $G=(V,E)$, where $V$ denotes the
set of nodes and $E$ denotes the set of edges, the goal is to
compute the correct answer $A$ that satisfies the requirements
specified in the task description.

Directly solving such problems using large language models
(LLMs) is challenging due to the difficulty of understanding
graph structures and the instability of direct code generation.
To address these challenges, PlanGraph decomposes the reasoning
process into three functional modules: \textit{planning generation},
\textit{plan-guided retrieval}, and \textit{execution-verified feedback reasoning}.

As illustrated in Figure~\ref{fig:framework}, the planning
module first converts the original task description into
structured pseudocode. The retrieval module augments the
reasoning process with external algorithmic knowledge.
Finally, the reasoning module generates executable programs
and verifies their correctness through iterative testing
and revision.

\subsection{Planning Module}

The planning module transforms an unstructured natural language
problem into a structured reasoning plan.

\subsubsection{Question Agent}

The question agent extracts key information from the original
task description in order to reduce the reasoning complexity
for downstream modules.

Formally, given the original question $Q$, the question agent
generates a structured representation of the task:

\begin{equation}
(T, G) \leftarrow QuestionAgent(Q)
\end{equation}

where $T$ denotes the task description and $G$ represents the
graph data extracted from the input.

The task description $T$ specifies the reasoning objective
and output requirements, while the graph data $G$ contains
the nodes, edges, and edge attributes used for graph
computation.

\subsubsection{Planning Agent}

After obtaining the structured task description $T$, the
planning agent generates a reasoning plan describing how
the problem should be solved.

\begin{equation}
P \leftarrow PlanningAgent(T)
\end{equation}

where $P$ denotes the generated planning pseudocode.

The reasoning plan can be viewed as a sequence of algorithmic
operations:

\begin{equation}
P = \{p_1, p_2, ..., p_n\}
\end{equation}

where each step $p_i$ corresponds to a specific graph
operation such as graph construction, traversal, or path
computation. The generated pseudocode serves as a structured
guide for downstream retrieval and reasoning modules.

\subsection{Plan-Guided Retrieval Module}

During complex problem solving, human experts often consult
external resources to obtain relevant algorithmic knowledge.
Inspired by this observation, PlanGraph introduces a retrieval
module that augments reasoning with information from graph
algorithm documentation.

\subsubsection{Knowledge Base}

We construct a knowledge base from documentation of graph
processing libraries such as NetworkX. Each entry in the
knowledge base corresponds to a specific API and contains
structured information describing the functionality of the
API, including its name, semantic description, input
parameters, return values, and example usage.

These documentation entries are converted into vector
representations and stored in a vector database to enable
efficient semantic retrieval.

\subsubsection{Search Agent}

Given the planning pseudocode $P$, the search agent retrieves
relevant knowledge entries from the documentation knowledge
base.

\begin{equation}
K \leftarrow SearchAgent(P)
\end{equation}

where $K$ denotes the retrieved knowledge set

\begin{equation}
K = \{k_1, k_2, ..., k_m\}
\end{equation}

and each element $k_i$ represents a relevant documentation
entry associated with a graph algorithm or API.

Compared with traditional question-based retrieval, the
plan-guided retrieval strategy leverages structured
pseudocode to produce more precise queries, and can further
refine or supplement retrieval when the returned knowledge
is insufficient for the current planning steps.

\subsection{Execution-Verified Feedback Reasoning Module}

The reasoning module performs graph computation based on the
planning pseudocode $P$, retrieved knowledge $K$, and graph
data $G$.

\subsubsection{Coding Agent}

The coding agent generates executable Python code using the
planning pseudocode and retrieved knowledge:

\begin{equation}
C_0 \leftarrow CodingAgent(P, K)
\end{equation}

where $C_0$ denotes the initial generated program. The
generated code typically invokes graph processing libraries
(e.g., NetworkX) to perform operations such as graph traversal,
shortest path computation, or combinational optimization.

\subsubsection{Execution and Verification}

To ensure reliability, PlanGraph introduces a verification
mechanism before applying the generated program to the
actual graph data.

First, the generated program is executed on a small test
graph to verify basic correctness. If the program fails
during testing, the system attempts to repair the code
using a reasoning agent.

The revision process can be expressed as

\begin{equation}
C_{t+1} \leftarrow ReasoningAgent(C_t, Error_t)
\end{equation}

where $C_t$ denotes the program generated at iteration $t$
and $Error_t$ represents the execution error observed during
testing.

If the revision process fails to produce a valid program
after several iterations, the system invokes the coding
agent again to generate a new candidate solution:

\begin{equation}
C_0' \leftarrow CodingAgent(P, K)
\end{equation}

This procedure forms a two-level feedback mechanism:
a \textit{repair loop} that attempts to fix incorrect
programs using execution feedback, and a \textit{regeneration loop}
that generates new programs when local repair fails or
important constraints are still missing.

Once a program successfully passes the unit test, it is
executed on the actual graph data $G$ to compute the final
result.

\subsubsection{Answer Generation}

After obtaining a verified program, the system executes it
on the input graph $G$ and obtains the computation result
$R$. The final answer is produced by formatting the result
according to the task requirements:

\begin{equation}
A \leftarrow FormatAnswer(R)
\end{equation}

The overall inference pipeline of PlanGraph is summarized
in Algorithm~\ref{alg:plangraph}.




\subsection{Design Rationale}

The design of PlanGraph is motivated by the observation that
directly applying large language models to graph reasoning
tasks often leads to unstable reasoning and hallucination.

When solving graph problems, the model must simultaneously
understand the problem description, recall relevant
algorithms, and generate executable code. These requirements
significantly increase the cognitive load of LLMs and often
result in incorrect reasoning steps.

To address this challenge, PlanGraph introduces an explicit
planning stage that separates high-level reasoning from
program execution. Instead of directly generating code,
the planning module first produces structured pseudocode
that decomposes the problem into a sequence of algorithmic
operations.

Another important component of the framework is the
plan-guided retrieval mechanism. The pseudocode
generated by the planning module contains explicit
algorithmic operations that can be directly mapped to
graph library APIs, enabling more accurate retrieval of
relevant algorithmic knowledge and clearer diagnosis of
retrieval failures.

Finally, PlanGraph incorporates an execution-feedback mechanism
to improve the reliability of generated solutions. By
introducing testing, repair, and regeneration loops, the
system can detect execution failures, localize errors, and
iteratively refine the generated program before applying it
to the target graph.

Together, these design choices enable PlanGraph to combine
the reasoning capabilities of LLMs with the precision of
graph algorithms, resulting in more reliable graph reasoning.

\begin{algorithm}[t]
\caption{PlanGraph Inference Pipeline}
\label{alg:plangraph}

\KwIn{Graph reasoning question $Q$}
\KwOut{Final answer $A$}

$(T,G) \leftarrow QuestionAgent(Q)$

$P \leftarrow PlanningAgent(T)$

$K \leftarrow SearchAgent(P)$

\While{solution not verified}{

Generate program $C$ using CodingAgent

Run unit test

\If{test fails}{
Revise program using ReasoningAgent
}

\If{revision fails}{
Regenerate program using CodingAgent
}

}

Execute $C$ on graph $G$

Extract result $R$

$A \leftarrow FormatAnswer(R)$

\Return $A$

\end{algorithm}

\subsection{Summary}

PlanGraph decomposes graph reasoning into three collaborative
stages: planning generation, plan-guided retrieval, and
execution-verified feedback reasoning. By combining
structured reasoning plans, retrieval grounded in planning
semantics, and feedback-driven program execution, the
framework enables reliable collaboration between language
models and graph computation tools.

\section{Experiments}

In this section, we evaluate the effectiveness of the proposed
\textbf{PlanGraph} framework on several graph reasoning benchmarks.
Our experiments aim to answer the following research questions:

\textbf{RQ1}: How does PlanGraph perform compared with existing
LLM-based graph reasoning systems?

\textbf{RQ2}: How does PlanGraph perform across different
graph reasoning benchmarks?

\textbf{RQ3}: How much improvement does PlanGraph achieve over
strong existing graph reasoning methods?

\textbf{RQ4}: How important are different modules in the framework?


\subsection{Experimental Setup}

\subsubsection{Datasets}

We evaluate PlanGraph on five widely used graph reasoning benchmarks:

\textbf{GraphArena}, \textbf{NLGraph}, \textbf{GraphWiz},
\textbf{LLM4DyG}, and \textbf{GraphInstruct}.

GraphArena focuses on classical graph computational problems,
including both polynomial-time tasks such as connectivity
analysis and shortest distance computation, as well as
NP-hard combinatorial optimization problems including
maximum clique, minimum vertex cover, and the traveling
salesman problem.

NLGraph is a natural language graph reasoning benchmark
that contains multiple algorithmic tasks such as connectivity,
maximum flow computation, cycle detection, topological sorting,
and shortest path queries.

GraphWiz evaluates graph reasoning ability through
instruction-style graph algorithm tasks and requires models
to follow algorithmic instructions to solve graph problems.

LLM4DyG focuses on reasoning over dynamic graphs where
temporal constraints are introduced into graph queries.
The benchmark evaluates whether models can correctly
perform reasoning over time-evolving graph structures.

GraphInstruct evaluates instruction-following capability
for graph reasoning tasks such as depth-first search,
neighbor queries, predecessor queries, and topological
sorting.

Unless otherwise specified, all reported results are evaluated
under the public benchmark settings and are computed using
exact-match comparison with ground-truth answers.


\subsubsection{Baselines}

We compare PlanGraph with three representative baselines.
First, we report the direct reasoning performance of
\textbf{GPT-4o}, which serves as a strong general-purpose LLM
baseline without task-specific graph reasoning scaffolding.
Second, we compare against \textbf{GraphTeam}, a representative
multi-agent graph reasoning framework. Third, we compare against
\textbf{GCoder}, a representative program-generation-based graph
reasoning method.

These two methods are chosen because they represent the two
strongest and most relevant technical routes for the evaluated
benchmarks: GraphTeam emphasizes multi-agent collaboration for
graph analysis, while GCoder emphasizes graph problem solving
through executable code generation. In the reported benchmark
results available from prior work, they also provide the strongest
or near-strongest published baselines.


\subsubsection{Evaluation Metric}

Following previous work on graph reasoning benchmarks,
we adopt \textbf{Exact Match (EM)} as the main evaluation metric.
A prediction is considered correct only if the generated
answer exactly matches the ground-truth result.


\begin{table*}[t]
\centering
\caption{Average exact match (\%) across graph reasoning benchmarks. 
Improvement is computed against the best baseline for each dataset.}

\setlength{\tabcolsep}{10pt}
\renewcommand{\arraystretch}{1.2}

\begin{tabular}{lcccccc}
\toprule
Dataset & \#Tasks & GPT-4o & GraphTeam & GCoder & PlanGraph & Improvement \\
\midrule
GraphArena & 10 & 49.08 & 95.14 & 94.27 & \textbf{97.74} & +2.60 \\
NLGraph & 8 & 51.40 & 97.71 & 96.90 & \textbf{98.72} & +1.01 \\
GraphWiz & 9 & 47.61 & 88.62 & 90.20 & \textbf{97.01} & +6.81 \\
LLM4DyG & 9 & 58.04 & \textbf{96.35} & 93.78 & 96.11 & -0.24 \\
GraphInstruct & 9 & 78.18 & 98.33 & 97.39 & \textbf{98.74} & +0.41 \\
\midrule
Average & -- & 56.86 & 95.23 & 94.51 & \textbf{97.66} & +2.12 \\
\bottomrule
\end{tabular}

\label{tab:overall}
\end{table*}


\subsection{Overall Performance}

Table~\ref{tab:overall} summarizes the average exact-match performance of
PlanGraph across five graph reasoning benchmarks. These
benchmarks cover a total of 45 task types included in the main evaluation, including
classical graph algorithms, combinatorial optimization
problems, instruction-following tasks, and dynamic graph
reasoning scenarios. The results in Table~\ref{tab:overall}
directly answer \textbf{RQ1} and \textbf{RQ3} by comparing
PlanGraph with existing LLM-based systems and measuring
its performance improvement over strong graph reasoning
baselines.

Overall, PlanGraph achieves the best performance on four
out of five benchmarks and improves the average exact-match
score from 95.23\% (best prior baseline) to 97.66\%, corresponding
to an average improvement of +2.12 percentage points.
These results show that the proposed planning-feedback
multi-agent framework consistently improves the reliability
of LLM-based graph reasoning systems.

The largest improvement is observed on the GraphWiz
benchmark, where PlanGraph increases the exact-match score from
90.20\% (GCoder) to 97.01\%, achieving a gain of
+6.81 points. GraphWiz contains several challenging
graph reasoning tasks such as Hamiltonian path detection
and subgraph matching, which require precise reasoning
over graph structures. The result indicates that the
planning-feedback reasoning pipeline can effectively guide
LLMs to perform structured reasoning, retrieve more relevant
algorithmic knowledge, and correct execution errors before
they propagate to final answers.

PlanGraph also achieves consistent improvements on
GraphArena and NLGraph, improving the best baseline
results by +2.60 and +1.01 points respectively.
These datasets mainly consist of classical graph
algorithm problems such as connectivity analysis,
maximum flow computation, and shortest path queries.
The results suggest that combining structured planning
with programmatic graph computation enables more reliable
execution of graph algorithms.

On the LLM4DyG benchmark, which focuses on temporal
graph reasoning tasks, PlanGraph achieves performance
comparable to the best existing method. Although the
result is slightly lower than the strongest prior baseline
(-0.24 points), the difference is marginal and the
framework still maintains strong performance across
dynamic graph queries.

Finally, PlanGraph achieves the highest exact-match score on the
GraphInstruct benchmark with 98.74\%, outperforming both
GraphTeam and GCoder. This result demonstrates that the
proposed framework generalizes well to instruction-based
graph reasoning tasks and can effectively handle a wide
range of graph analysis problems.

Overall, these results show that the planning-feedback
multi-agent reasoning framework consistently improves
the performance of LLM-based graph reasoning systems
across diverse benchmarks.


\subsection{Analysis by Task Categories}

To further analyze the reasoning capabilities of PlanGraph,
we examine its performance across different categories of
graph reasoning tasks.

Specifically, we group tasks from multiple benchmarks into
several representative categories, including graph structure
analysis, shortest path computation, matching problems,
flow optimization, and Hamiltonian problems. These categories
cover a broad range of classical graph algorithms and
combinatorial reasoning tasks.

\begin{figure}[h]
\centering
\includegraphics[width=\columnwidth]{task_category.pdf}
\caption{Performance comparison across graph reasoning task categories.}
\label{fig:task_category}
\end{figure}

Figure~\ref{fig:task_category} presents the comparison
between GPT-4o, strong prior methods, and
PlanGraph across these task categories.

As shown in Figure~\ref{fig:task_category}, PlanGraph achieves
the highest exact-match performance on matching problems and graph structure
analysis tasks, reaching nearly perfect performance.
These problems can typically be solved using well-defined
graph algorithms such as bipartite matching, connectivity
analysis, and graph traversal, which can be effectively
represented through the generated planning pseudocode.

For shortest path and Hamiltonian problems, PlanGraph also
achieves very high performance, demonstrating that the generated
programs can reliably perform graph traversal and
combinatorial search operations. The planning stage provides
a structured reasoning process that allows the system to
identify the appropriate algorithms and execute them through
graph computation libraries.

Flow optimization tasks show slightly lower performance compared
with other categories. This is mainly because flow algorithms
often involve more complex implementation details such as
residual graphs and capacity constraints, which increases
the difficulty of code generation and reasoning.

Overall, the results indicate that PlanGraph performs
robustly across diverse graph reasoning tasks and generalizes
well across multiple benchmarks.

\subsection{Case Study}

To further understand how \textbf{PlanGraph} solves graph reasoning
problems, we present a qualitative case study on a Traveling
Salesman Problem (TSP) instance from the GraphArena benchmark.
The purpose of this analysis is to illustrate the internal
reasoning process of the proposed framework and demonstrate
how the planning-feedback pipeline improves reasoning reliability.

\begin{figure}[t]
\centering
\includegraphics[width=\columnwidth]{case1.pdf}
\caption{Example TSP instance used in the case study.}
\label{fig:case_tsp}
\end{figure}


Figure~\ref{fig:case_tsp} illustrates a representative TSP
problem used in the case study. The graph contains five nodes
with weighted edges representing pairwise distances between
nodes. The objective is to find the shortest Hamiltonian cycle
starting from node A that visits each node exactly once and
returns to the starting node.

Directly solving such combinatorial optimization problems
with large language models is challenging because the model
must simultaneously understand the graph structure, identify
the appropriate algorithm, and generate executable code.
Without structured reasoning guidance, direct prompting
methods often produce incomplete paths or violate the
Hamiltonian constraint.

To address this issue, PlanGraph decomposes the reasoning
process into structured stages including planning generation,
plan-guided retrieval, and execution-verified feedback reasoning.

\begin{figure}[t]
\centering
\includegraphics[width=\columnwidth]{case2.pdf}
\caption{Reasoning process of PlanGraph on the TSP example.}
\label{fig:case_pipeline}
\end{figure}

Figure~\ref{fig:case_pipeline} illustrates the reasoning
pipeline of PlanGraph when solving the TSP instance. First,
the planning module analyzes the task description and
generates structured pseudocode that decomposes the problem
into algorithmic operations such as graph construction,
candidate tour enumeration, and path cost computation.
This planning step provides a clear reasoning structure
for the downstream modules and makes key constraints explicit
before program generation starts.

Next, the retrieval module identifies relevant algorithmic
knowledge from graph API documentation based on the
generated pseudocode. When the initial evidence is not
sufficient, the system can refine the retrieved snippets
or supplement them with additional retrieval conditioned on
the current planning steps and missing constraints.

Using the planning pseudocode and retrieved knowledge,
the coding agent generates executable Python code that
implements the corresponding graph computation. Before
applying the generated program to the target graph,
PlanGraph performs execution verification by running the
program on a test graph. If execution errors occur, the
reasoning agent revises the generated code using execution
feedback. If local repair is insufficient, the system can
regenerate a new candidate under the same planning guidance.

Finally, the verified program is executed on the input
graph to compute the optimal tour. In this example, the
system returns the path
\[
A \rightarrow B \rightarrow E \rightarrow D \rightarrow C \rightarrow A
\]
with a total cost of 14.

This case study demonstrates how PlanGraph decomposes a
complex graph reasoning task into structured planning,
knowledge retrieval, and feedback-driven program execution.
By separating planning from implementation and introducing
execution feedback, the framework reduces hallucination,
improves constraint satisfaction, and strengthens the
reliability of LLM-based graph reasoning systems.

\subsection{Ablation Study}

To evaluate the contribution of each component in the
PlanGraph framework, we conduct an ablation study by
removing individual modules from the system.

Specifically, we evaluate three variants:

\textbf{w/o Planning}: removing the planning module and
directly generating code from the question.

\textbf{w/o Retrieval}: removing the plan-guided
knowledge retrieval module.

\textbf{w/o Verification}: removing the execution
verification and code revision mechanism.

\begin{table}[t]
\centering
\caption{Ablation study of PlanGraph.}
\label{tab:ablation}
\begin{tabular}{lccc}
\toprule
Method & GraphArena & NLGraph & GraphWiz \\
\midrule
PlanGraph & \textbf{97.88} & \textbf{98.72} & \textbf{97.01} \\
w/o Planning & 86.35 & 89.14 & 84.27 \\
w/o Retrieval & 94.12 & 95.33 & 92.48 \\
w/o Verification & 92.67 & 94.01 & 90.72 \\
\bottomrule
\end{tabular}
\end{table}

Table~\ref{tab:ablation} presents the ablation results on
three representative benchmarks. The results clearly
demonstrate the contribution of each component in the
PlanGraph framework.

Removing the planning module leads to the largest
performance drop across all datasets. On GraphArena,
the exact-match score decreases from 97.74\% to 86.35\%.
This result confirms that the planning stage plays a
critical role in guiding the reasoning process.
Without structured pseudocode, the model must directly
generate executable programs from natural language
descriptions, which significantly increases the
difficulty of algorithm selection and reasoning.

Removing the verification mechanism also causes
noticeable performance degradation. The exact-match score drops
to 92.67\% on GraphArena and 90.72\% on GraphWiz.
This observation indicates that execution feedback is
essential for detecting and correcting errors in
generated programs. The verification loop enables the
system to iteratively repair incorrect code before
applying it to the target graph.

The retrieval module provides additional algorithmic
knowledge by linking the generated planning pseudocode
to relevant graph APIs. When this component is removed,
the performance decreases moderately but remains higher
than the variants without planning or verification.
This suggests that while retrieval improves reasoning
efficiency and code generation quality, the planning
and verification modules play more critical roles in
the overall reasoning pipeline.

Overall, these results validate the design of the
PlanGraph framework and demonstrate that planning,
retrieval, and feedback-driven verification work
together to improve the reliability of LLM-based
graph reasoning.

\section{Conclusion}

In this paper, we propose \textbf{PlanGraph}, a planning-feedback multi-agent framework for improving large language model reasoning on graph computation tasks. By introducing an explicit planning stage, PlanGraph decomposes complex graph problems into structured algorithmic steps and uses plan-guided retrieval to support reliable program generation. The framework further integrates execution-based feedback to refine generated programs, correct omitted constraints, and reduce hallucination errors. Extensive experiments on multiple graph reasoning benchmarks show that PlanGraph achieves the best performance on four of the five evaluated benchmarks, reaches an average exact-match score of 97.66\%, and improves overall performance over strong existing methods.

\nocite{*}
\bibliographystyle{ACM-Reference-Format}
\bibliography{references}


\appendix

\section{Detailed Benchmark Results}

Due to space limitations in the main paper, this appendix
reports the detailed per-task results of PlanGraph on all
evaluated benchmarks. For completeness, we present the
performance of PlanGraph together with several representative
baselines on each individual task. All reported results are
reported under the corresponding public evaluation settings. Unless otherwise
specified, the base language model used in our framework is
GPT-4o.

The reported benchmarks cover a diverse set of graph reasoning
tasks, including classical graph algorithms, combinatorial
optimization problems, and instruction-based graph queries,
providing a comprehensive view of the performance of PlanGraph
across different graph reasoning scenarios.

\subsection{GraphArena}
\vspace{-6pt}
\begin{table}[H]
\centering
\caption{Detailed results on GraphArena (\%).}
\begin{tabular}{lcccc}
\toprule
Task & GPT-4o & GraphTeam & GCoder & PlanGraph \\
\midrule
CN & 100.0 & 100.0 & 100.0 & 100.0 \\
CC & 100.0 & 100.0 & 100.0 & 100.0 \\
SD & 100.0 & 100.0 & 100.0 & 100.0 \\
GD & 100.0 & 100.0 & 100.0 & 100.0 \\
GED & 23.3 & 91.0 & 89.5 & 96.5 \\
MIS & 22.4 & 97.8 & 96.8 & 98.8 \\
MVC & 19.3 & 98.6 & 97.6 & 100.0 \\
MCP & 12.5 & 90.5 & 88.7 & 100.0 \\
MCS & 8.4 & 90.8 & 89.4 & 95.3 \\
TSP & 4.9 & 82.7 & 80.7 & 86.8 \\
\bottomrule
\end{tabular}
\end{table}

\subsection{NLGraph}
\vspace{-6pt}
\begin{table}[H]
\centering
\caption{Detailed results on NLGraph (\%).}
\begin{tabular}{lcccc}
\toprule
Task & GPT-4o & GraphTeam & GCoder & PlanGraph \\
\midrule
CONN & 92.10 & 100.0 & 100.0 & 100.0 \\
MF & 67.30 & 93.67 & 94.10 & 97.53 \\
CYC & 95.80 & 100.0 & 100.0 & 100.0 \\
TOPO & 60.50 & 91.20 & 92.80 & 94.80 \\
SP & 72.40 & 100.0 & 100.0 & 100.0 \\
BM & 81.30 & 100.0 & 100.0 & 100.0 \\
HP & 40.70 & 100.0 & 100.0 & 100.0 \\
GNN & 52.10 & 96.80 & 97.40 & 97.40 \\
\bottomrule
\end{tabular}
\end{table}

\subsection{GraphWiz}
\begin{table}[H]
\centering
\caption{Detailed results on GraphWiz (\%).}
\begin{tabular}{lcccc}
\toprule
Task & GPT-4o & GraphTeam & GCoder & PlanGraph \\
\midrule
CYC & 95.00 & 100.0 & 100.0 & 100.0 \\
CONN & 70.30 & 86.50 & 90.20 & 91.20 \\
BIP & 82.70 & 96.30 & 98.00 & 100.0 \\
TOPO & 75.20 & 100.0 & 100.0 & 96.50 \\
SP & 78.60 & 95.00 & 96.40 & 98.50 \\
MTS & 71.80 & 94.00 & 95.10 & 93.80 \\
MF & 69.10 & 94.00 & 94.70 & 95.10 \\
HP & 15.40 & 32.50 & 40.20 & 99.20 \\
SM & 88.30 & 99.30 & 99.50 & 98.80 \\
\bottomrule
\end{tabular}
\end{table}

\subsection{LLM4DyG}

\begin{table}[H]
\centering
\small
\caption{Detailed results on LLM4DyG (\%).}
\begin{tabular}{lccccc}
\toprule
Task & GPT-4o & LLM4DyG & GraphTeam & GCoder & PlanGraph \\
\midrule
WL & 58.00 & 47.30 & 96.33 & 94.33 & 98.67 \\
WC & 67.33 & 64.10 & 94.33 & 94.78 & 93.56 \\
WTC & 87.67 & 63.00 & 100.00 & 97.08 & 100.00 \\
WN & 33.33 & 45.90 & 95.33 & 94.74 & 97.47 \\
WNB & 44.67 & 20.00 & 100.00 & 94.83 & 100.00 \\
CTC & 77.33 & 66.70 & 99.67 & 100.00 & 99.33 \\
CTP & 57.67 & 59.80 & 87.00 & 100.00 & 92.03 \\
FTP & 46.67 & 80.80 & 81.00 & 74.58 & 84.57 \\
SE & 49.67 & 35.80 & 99.67 & 93.67 & 99.33 \\
\bottomrule
\end{tabular}
\end{table}

\subsection{GraphInstruct}

\begin{table}[H]
\centering
\small
\caption{Detailed results on GraphInstruct (\%).}
\begin{tabular}{lccccc}
\toprule
Task & GPT-4o & GraphInstruct & GraphTeam & GCoder & PlanGraph \\
\midrule
NB & 98.00 & 99.00 & 100.00 & 99.53 & 100.00 \\
BIP & 89.67 & 85.00 & 99.00 & 100.00 & 100.00 \\
EDGE & 90.67 & 84.00 & 99.33 & 97.47 & 98.63 \\
CONN & 92.22 & 83.00 & 100.00 & 100.00 & 100.00 \\
DEG & 98.00 & 81.00 & 100.00 & 98.67 & 100.00 \\
DFS & 89.33 & 46.00 & 95.67 & 94.33 & 96.33 \\
PRED & 47.41 & 41.00 & 98.00 & 99.75 & 99.75 \\
TOPO & 59.00 & 33.00 & 98.33 & 99.50 & 99.25 \\
CC & 39.30 & 83.00 & 94.67 & 83.30 & 94.74 \\
\bottomrule
\end{tabular}
\end{table}


\end{document}
